\chapter{Introducción}
\pagestyle{fancy}

\section{Descripción de la Misión}
INTISAT es una misión de demostración tecnológica en formato \textbf{CubeSat 2U} liderada por la Pontificia Universidad Católica del Perú (PUCP) a través del Instituto de Radioastronomía (INRAS). El objetivo principal es validar la integración de una carga útil compacta de microscopía por imagen para operación en órbita terrestre, desarrollada por estudiantes e investigadores como parte de los esfuerzos de PUCP en educación y desarrollo de tecnología espacial.

INTISAT demostrará una plataforma totalmente integrada compuesta por una computadora de a bordo (OBC), un sistema de energía eléctrica (EPS), un sistema de telemetría, rastreo y comando (TTC), componentes estructurales y una carga útil científica. Todo el sistema está contenido en una envolvente mecánica CubeSat 2U. La implementación actual corresponde a un \textit{Modelo de Ingeniería}, base para mejoras posteriores y una eventual calificación de vuelo.

\insertfigure{IntiSat_Renderization.png}{0.25}{Vista del modelo 3D de INTISAT.}

\section{Objetivos de la Misión}
Los objetivos principales de INTISAT son:
\begin{itemize}
    \item Validar la integración y operación de un sensor de imagen miniaturizado para microscopía en órbita terrestre baja (LEO).
    \item Demostrar la funcionalidad de subsistemas desarrollados por estudiantes operando dentro de una plataforma de nanosatélite unificada.
    \item Desarrollar y verificar una arquitectura CubeSat escalable y modular alineada con estándares académicos y educativos.
    \item Servir como plataforma de entrenamiento práctico en ingeniería de sistemas, diseño, integración y pruebas de satélites.
\end{itemize}

\section{Concepto de Operaciones (ConOps)}
Siguiendo enfoques estandarizados de operación satelital, el ciclo operativo de INTISAT se define mediante cuatro modos de misión claramente delimitados. La gestión y transiciones entre modos las ejecuta la OBC (ver Cap.~3).

\begin{itemize}
    \item \textbf{Modo Despliegue e Inicialización:} durante el lanzamiento todos los sistemas están físicamente deshabilitados por \textit{kill-switches}. Tras la eyección, los sistemas se energizan, la OBC arranca, se ejecutan verificaciones iniciales de \textit{housekeeping} y se inicia la secuencia de despliegue de antenas después de una ventana de silencio obligatoria de 40~minutos.

    \item \textbf{Modo Nominal (Rutina):} estado operativo estándar en el que el EPS maximiza la recolección solar, la baliza transmite periódicamente y la carga útil (Microscopía / IA) se activa durante ventanas predeterminadas sobre el territorio peruano u otros objetivos planificados.

    \item \textbf{Modo Contingencia (Seguridad):} estado seguro al que se entra automáticamente si se detecta una anomalía crítica (por ejemplo, SoC bajo de batería o eventos \textit{latch-up}). Las cargas no esenciales se deshabilitan hasta la intervención del control terreno.

    \item \textbf{Modo Decomisión:} estado de fin de vida activado mediante telecomando desde tierra o por temporizador programado. Asegura la pasivación segura de la batería y el apagado de transceptores antes del reingreso atmosférico.
\end{itemize}

\section{Miembros del Proyecto}
INTISAT es una iniciativa académica impulsada por equipos interdisciplinarios de PUCP. La estructura del proyecto incluye:
\begin{itemize}
    \item \textbf{INRAS – PUCP:} coordinación de misión, ingeniería de sistemas, desarrollo del módulo TTC e integración de la plataforma.
    \item \textbf{Facultad de Ciencias e Ingeniería – PUCP:} desarrollo de software, electrónica y subsistemas mecánicos.
    \item \textbf{Socios externos (fase futura):} colaboradores para pruebas, coordinación de lanzamiento u operaciones de misión.
\end{itemize}

\section{Logotipo de Misión}
\insertfigure{pucp_inras_logo.png}{0.25}{Logotipo oficial de la misión INTISAT.}
